\documentclass[11pt]{article}
\usepackage{icmlko}
\begin{document}

\title{판이 낮은 것은 모음 탓이 아니라 시간 탓이다 --- 천장을 재려다
내 물음이 틀렸다는 걸 알았다}
\author{월드모델 실험실}
\date{노트 445 · 2026년 8월 2일}
\maketitle

\begin{abstract}
노트 444 가 집 밖 짝에서 한 일을 판에 그대로 걸어 봤다. 도메인마다
\emph{자기 행만으로} 교차검증 천장을 재고 판의 유보 rho 를 그 천장으로
나눴다. 첫 표는 그럴듯했다 --- 만화가 비 $0.55$ 로 꼴찌이니 \textbf{제일
큰 도메인이 모음에게 제일 홀대받는다}는 이야기가 되는 듯했다.

\textbf{그 이야기는 틀렸다.} 교차검증 천장은 \emph{같은 시기} 행끼리 서로
맞히는 것이고 판은 \emph{2025년 이후}를 맞히는 것이다. 둘을 바로 견주면
\textbf{앞을 내다보는 비용이 모음 탓으로 잘못 찍힌다}. 그래서 같은 자를
하나 더 놨다 --- \textbf{도메인이 자기 행만으로, 판과 똑같은 시간
분할로} 예보하게 한 것이다. 그러면 값이 둘로 갈린다.

\begin{center}
\begin{tabular}{lrr}
\hline
 & 평균 & 부호 \\
\hline
\textbf{예보 값} (혼자 $-$ 같은 시기 천장) & $\mathbf{-0.1492}$ & 음수 9/10 \\
모음 값 (판 $-$ 혼자) & $+0.0699$ & 양수 8/10 \\
\hline
\end{tabular}
\end{center}

\textbf{열 도메인 중 일곱에서 예보 값이 모음 값보다 크다.} 만화는 비
$0.55$ 로 홀대받는 게 아니었다 --- 같은 시기 천장 $+0.6516$ 에서 혼자
예보하면 $+0.3388$ 로 떨어지고($-0.3128$), 판은 거기에 $+0.0217$ 을
\emph{보탠다}. 병은 모음이 아니라 \textbf{시간}이다.
\end{abstract}

\begin{figure}[t]
\centering
\includegraphics[width=\textwidth]{figs/f1.pdf}
\caption{\textbf{가} --- 도메인마다 검은 동그라미(같은 시기 천장)에서
붉은 막대만큼 떨어져 혼자 예보에 이르고, 초록 막대만큼 올라 판에 이른다.
붉은 막대가 거의 어디서나 길다. \textbf{나} --- 모음 값은 학습 행 수와
반대로 간다(rho $-0.539$). 도서 80행에 $+0.2551$, 만화 5{,}647행에
$+0.0217$.}
\end{figure}

\section{모음은 작은 도메인을 위한 것이다}

갈라 놓고 보니 모음 값의 정체가 분명하다. \textbf{학습 행 수와 반대로
간다}(스피어만 $-0.539$).

\begin{center}
\begin{tabular}{lrr}
\hline
 & 모음 값 & 도메인 \\
\hline
작은 도메인 (학습 $<300$) & $\mathbf{+0.1384}$ & 도서 · 아이돌 · 게임 · 시장팝업 \\
큰 도메인 (학습 $\ge 1500$) & $+0.0242$ & 만화 · 애니 · 세계애니 · 웹툰 · 모바일 · 펀딩 \\
\hline
\end{tabular}
\end{center}

도서는 학습 80행으로 혼자서는 $+0.1177$ 밖에 못 내는데 판에 붙으면
$+0.3728$ 이다. 반대쪽 끝에서 모바일($-0.0130$)과 애니($-0.0102$)는
\textbf{판에 붙어서 아주 조금 손해}를 본다 --- 둘 다 잡음 크기지만
부호는 정직하게 적는다.

\section{시장팝업 --- 자기 자료에는 아무것도 없다}

한 도메인만 반대로 간다. \textbf{시장팝업의 같은 시기 천장이 $-0.0901$
이다} --- 자기 축으로 자기 라벨을 교차검증하면 $0$ 보다 못하다. 그런데
혼자 예보는 $+0.0113$, 판에 붙으면 $+0.1340$ 이다.

\textbf{시장팝업이 가진 $+0.134$ 는 전부 빌린 것이다.} 판에서 제일 낮은
점수인데, 낮은 이유가 ``덜 배웠다''가 아니라 \textbf{자기 자료에 든 게
없다}는 것이다. 노트 412 가 ``시장팝업 전용 축은 그 도메인 $+0.0241$ 인데
판엔 못 온다''고 했던 것의 뒷면이다 --- 전용 축이 그 도메인을 조금
올리는 것도, 그 도메인의 바닥이 $0$ 아래라서다.

\section*{규약 --- 여지를 말하려면 시간 분할을 맞추고 말한다}

이 노트가 만드는 규약. \textbf{도메인의 ``남은 여지''를 말할 때는 천장과
점수의 \emph{시간 분할을 맞춘 뒤} 말한다.} 같은 시기 교차검증 천장과
앞을 내다본 유보 rho 를 바로 나누면, \textbf{예보 비용이 통째로 모형
탓으로 찍힌다} --- 이 노트의 첫 표가 정확히 그 실수였고, 그대로 뒀으면
``만화를 모음에서 빼자''는 엉뚱한 갈래로 갔을 것이다.

\section*{남은 것}

\begin{itemize}
\item \textbf{예보 값이 진짜 병이면 시간 축을 봐야 한다.} 지금 판에
들어간 시간 관련 축은 달력 넷과 추세뿐이다. 노트 445 까지 시간은
\emph{채점 방식}이었지 \emph{모형의 재료}가 아니었다.
\item \textbf{시장팝업의 천장이 왜 음수인가.} 태거 갈래다 --- CN 과 같은
병일 수 있다(노트 444).
\item \textbf{모바일 · 애니는 판에서 빼면 오르나.} $-0.013$ · $-0.010$ 은
잡음 크기라 그 자체로는 못 정한다. 관문에 올리려면 짝 12뽑기가 필요하다.
\end{itemize}

\end{document}
